\chapter*{符号说明}


我在本书中采用了以下常规记号。


\subsection*{数学}

\begin{tabular}{ll}
\hline 
$\binom{n}{m}$ & “从$n$个元素中选出$m$个”的组合数，其值为$\frac{n!}{m!(n-m)!}$。 \\
$\sum$ & 求和运算，例如：$\sumn a_i = a_1 + \cdots + a_n$ \\
$I(\cdot)$ & 指示函数，即：若事件$A$发生则$I(A)=1$，否则为0。\\
$\#$ & 计算集合中单位的数量\\
$\approx$ & 大致相等 \\
$\propto$ & 与某值成正比（忽略一些不重要的常数后）。 \\
$\text{logit}$ &  $\text{logit}(x) = \log \frac{x}{1-x}$ \\
$\expit$ & $\expit(x) =   \frac{e^x}{1+e^x} =  (1 + e^{-x})^{-1}$ \\
\hline 
\end{tabular}

\subsection*{基本概率与统计学}

\begin{tabular}{ll}
\hline 
$\pr(\cdot)$ & 概率 \\
$E(\cdot)$ & 随机变量的期望值\\
$\var(\cdot)$ & 随机变量的方差 \\
$\cov(\cdot)$ & 随机变量之间的协方差\\
$\asim$ & “$A \asim B$”表示$A$与$B$具有相同的渐近分布。\\
$\rho_{YX}$ & $Y$与$X$之间的皮尔逊相关系数 \\
$R^2_{YX}$ & $Y$与$X$之间的平方多重相关系数 \\
\hline 
\end{tabular}


\subsection*{随机变量}

\begin{tabular}{ll}
\hline 
伯努利分布$(p)$ & 概率为 $p$ 的伯努利分布\\
二项分布$(n,p)$& 进行$n$次试验且成功概率为$p$的二项分布\\
$\textsc{N}(\mu, \sigma^2)$ & 均值为$\mu$、方差为$\sigma^2$的正态分布\\
$t_\nu$ & 自由度为$\nu$的$t$分布 \\
$\chi^2_\nu$ & 自由度为$\nu$的卡方分布\\
$z_{1-\alpha/2}$ & $\N01$分布的$1-\alpha/2$分位数，例如$z_{0.975}=1.96$。\\
\hline 
\end{tabular}


\subsection*{因果推断}

\begin{tabular}{ll}
\hline 
 $Y_i(1), Y_i(0)$ & 单位$i$在处理与对照条件下的潜在结果\\
$\tau$ & 有限总体平均因果效应 $\tau = n^{-1}\sumn \{Y_i(1)-Y_i(0)\}$ \\
$\tau$ & 超总体平均因果效应 $\tau = E \{Y(1)-Y(0)\}$\\
$\mu_1(X)$ & 结果模型 $\mu_1(X) = E(Y\mid Z=1, X)$\\
$\mu_0(X)$ & 结果模型 $\mu_0(X) = E(Y\mid Z=0, X)$\\
$e(X)$ & 倾向得分 $e(X) = \pr(Z=1\mid X)$ \\
$\tau_\cp$ & 依从者的平均因果效应 $ \tau_\cp = E\{ Y(1) - Y(0)\mid U=\cp \} $ \\
$U$ & 未测量混杂因素 \\
$U$ & 潜在的依从状态 $U=(D(1), D(0))$ \\
$Y(z,M_{z'})$ & 嵌套式潜在结果（用于中介分析）\\
 \hline 
\end{tabular}